% Chunk: Chapter 7 introduction and Section 7.1, Canonical Variables
% Source: Weinberg QFT Vol. I, printed pp. 292--298 (physical PDF pp. 315--321)
% Status: transcribed; source-reviewed; compile-clean; render-reviewed
% Physical pages: 315--321 (through the end of Section 7.1). Figures: none. Uncertainties: none.

Ever since the birth of quantum field theory in the papers of Born, Dirac,
Fermi, Heisenberg, Jordan, and Pauli in the late 1920s, its development
has been historically linked to the canonical formalism, so much so that it
seems natural to begin any treatment of the subject today by postulating a
Lagrangian and applying to it the rules of canonical quantization. This is
the approach used in most books on quantum field theory. Yet historical
precedent is not a very convincing reason for using this formalism. If
we discovered a quantum field theory that led to a physically satisfactory
$S$-matrix, would it bother us if it could not be derived by the canonical
quantization of some Lagrangian?

To some extent this question is moot because, as we shall see in Section
7.1, all of the most familiar quantum field theories furnish canonical
systems, and these can easily be put in a Lagrangian form. However,
there is no proof that every conceivable quantum field theory can be
formulated in this way. And even if it can, this does not in itself explain
why we should \textit{prefer} to use the Lagrangian formalism as a starting point
in constructing various quantum field theories.

The point of the Lagrangian formalism is that it makes it easy to
satisfy Lorentz invariance and other symmetries: a classical theory with
a Lorentz-invariant Lagrangian density will when canonically quantized
lead to a Lorentz-invariant quantum theory. That is, we shall see here
that such a theory allows the construction of suitable quantum mechanical
operators that satisfy the commutation relations of the Poincar\'e algebra,
and therefore leads to a Lorentz-invariant $S$-matrix.

This is not so trivial. We saw in the previous chapter that in theories
with derivative couplings or spins $j\geq1$, it is not enough to take the
interaction Hamiltonian as the integral over space of a scalar interaction
density; we also need to add non-scalar terms to the interaction density
to compensate for non-covariant terms in the propagators. The canonical
formalism with a scalar Lagrangian density will automatically provide
these extra terms. Later, when we come to non-Abelian gauge theories
in Volume II, this extra convenience will become a necessity; it would be
just about hopeless to try to guess at the form of the Hamiltonian in such
theories without starting with a Lorentz-invariant and gauge-invariant
Lagrangian density.

\subsection{Canonical Variables}

In this section we shall show that various quantum field theories that we
have constructed so far satisfy the commutation rules and equations of
motion of the Hamiltonian version of the canonical formalism. It is the
Hamiltonian formalism that is needed to calculate the $S$-matrix (whether
by operator or path-integral methods) but it is not always easy to choose
Hamiltonians that yield a Lorentz-invariant $S$-matrix. In the balance
of this chapter we shall take the Lagrangian version of the canonical
formalism as our starting point, and use it to derive physically satisfactory
Hamiltonians. The purpose of the present section is to identify the
canonical fields and their conjugates in various field theories, to tell us
how to separate the free-field terms in the Lagrangian, and incidentally to
reassure us that the canonical formalism is indeed applicable to physically
realistic theories.

We first show that the free fields constructed in Chapter 5 automatically
provide a system of quantum operators $q^n(\mathbf{x},t)$ and canonical conjugates
$p_n(\mathbf{x},t)$ that satisfy the familiar canonical commutation or anticommutation
relations:
\begin{equation}
\begin{aligned}
[q^n(\mathbf{x},t),p_{\bar n}(\mathbf{y},t)]
&=i\delta^3(\mathbf{x}-\mathbf{y})\delta^n_{\bar n},
&&\text{if either particle is a boson},\\
\{q^n(\mathbf{x},t),p_{\bar n}(\mathbf{y},t)\}
&=i\delta^3(\mathbf{x}-\mathbf{y})\delta^n_{\bar n},
&&\text{if both particles are fermions}.
\end{aligned}
\tag{7.1.1}\label{eq:7.1.1}
\end{equation}
\begin{equation}
\begin{aligned}
[q^n(\mathbf{x},t),q^{\bar n}(\mathbf{y},t)]&=0,
&&\text{if either particle is a boson},\\
\{q^n(\mathbf{x},t),q^{\bar n}(\mathbf{y},t)\}&=0,
&&\text{if both particles are fermions}.
\end{aligned}
\tag{7.1.2}\label{eq:7.1.2}
\end{equation}
\begin{equation}
\begin{aligned}
[p_n(\mathbf{x},t),p_{\bar n}(\mathbf{y},t)]&=0,
&&\text{if either particle is a boson},\\
\{p_n(\mathbf{x},t),p_{\bar n}(\mathbf{y},t)\}&=0,
&&\text{if both particles are fermions}.
\end{aligned}
\tag{7.1.3}\label{eq:7.1.3}
\end{equation}
In each equation, the commutator line applies if either of the particles
created and destroyed by the two operators are bosons, and the
anticommutator line applies if both particles are fermions. For instance, the real
scalar field $\phi(x)$ for a self-charge-conjugate particle of zero spin was
found in Section 5.2 to obey the commutation relation
\[
[\phi(x),\phi(y)]=\Delta(x-y),
\]
where $\Delta$ is the function
\[
\Delta(x-y)\equiv\int\frac{d^3k}{2k^0(2\pi)^3}
\left[e^{ik\cdot(x-y)}-e^{-ik\cdot(x-y)}\right]
\]
with $k^0=\sqrt{\mathbf{k}^2+m^2}$. We note that
\[
\Delta(\mathbf{x},0)=0,
\qquad
\dot\Delta(\mathbf{x},0)=-i\delta^3(\mathbf{x}).
\]
(A dot denotes the derivative with respect to the time $x^0$.) It is easy
then to see that the field and its time-derivative $\dot\phi$ obey the equal-time
commutation relations:
\begin{equation}
[\phi(\mathbf{x},t),\dot\phi(\mathbf{y},t)]
=i\delta^3(\mathbf{x}-\mathbf{y}),
\tag{7.1.4}\label{eq:7.1.4}
\end{equation}
\begin{equation}
[\phi(\mathbf{x},t),\phi(\mathbf{y},t)]=0,
\tag{7.1.5}\label{eq:7.1.5}
\end{equation}
\begin{equation}
[\dot\phi(\mathbf{x},t),\dot\phi(\mathbf{y},t)]=0.
\tag{7.1.6}\label{eq:7.1.6}
\end{equation}
Therefore we may define canonical variables
\begin{equation}
q(\mathbf{x},t)\equiv\phi(\mathbf{x},t),
\qquad
p(\mathbf{x},t)\equiv\dot\phi(\mathbf{x},t)
\tag{7.1.7}\label{eq:7.1.7}
\end{equation}
which satisfy the canonical commutation relations~\eqref{eq:7.1.1}--\eqref{eq:7.1.3}.

For the complex scalar field of a particle of spin zero with a distinct
antiparticle, the commutation relations are
\[
[\phi(x),\phi^\dagger(y)]=\Delta(x-y),
\qquad
[\phi(x),\phi(y)]=0.
\]
We may therefore define the free-particle canonical variables as the complex
operators
\begin{equation}
q(\mathbf{x},t)\equiv\phi(\mathbf{x},t),
\tag{7.1.8}\label{eq:7.1.8}
\end{equation}
\begin{equation}
p(\mathbf{x},t)\equiv\dot\phi^\dagger(\mathbf{x},t).
\tag{7.1.9}\label{eq:7.1.9}
\end{equation}
Equivalently, writing $\phi\equiv(\phi_1+i\phi_2)/\sqrt{2}$ with $\phi_k$ Hermitian for
$k=1,2$, we have canonical variables
\begin{equation}
q^k(\mathbf{x},t)=\phi_k(\mathbf{x},t),
\tag{7.1.10}\label{eq:7.1.10}
\end{equation}
\begin{equation}
p_k(\mathbf{x},t)=\dot\phi_k(\mathbf{x},t),
\tag{7.1.11}\label{eq:7.1.11}
\end{equation}
and these satisfy the commutation relations~\eqref{eq:7.1.1}--\eqref{eq:7.1.3}.

For the real vector field of a particle of spin one, the commutation
relations are given by Section 5.3 as
\[
[v^\mu(x),v^\nu(y)]
=\left[\eta^{\mu\nu}-\frac{\partial^\mu\partial^\nu}{m^2}\right]\Delta(x-y).
\]
(We are using $v^\mu$ rather than $V^\mu$ for the vector field because we want to
reserve upper case letters here for the fields in the Heisenberg picture.)
Here the free-particle canonical variables may be taken as
\begin{equation}
q^i(\mathbf{x},t)=v^i(\mathbf{x},t),
\tag{7.1.12}\label{eq:7.1.12}
\end{equation}
\begin{equation}
p_i(\mathbf{x},t)=\dot v^i(\mathbf{x},t)+\frac{\partial v^0(\mathbf{x},t)}{\partial x^i},
\tag{7.1.13}\label{eq:7.1.13}
\end{equation}
with $i=1,2,3$. The reader may check that \eqref{eq:7.1.12} and \eqref{eq:7.1.13} satisfy
the commutation relations~\eqref{eq:7.1.1}--\eqref{eq:7.1.3}. The field equations~\eqref{eq:5.3.36} and
\eqref{eq:5.3.38} together with Eq.~\eqref{eq:7.1.13} allow us to express $v^0$ in terms of the
other variables as
\begin{equation}
v^0=m^{-2}\boldsymbol{\nabla}\cdot\mathbf{p},
\tag{7.1.14}\label{eq:7.1.14}
\end{equation}
so $v^0$ is not regarded as one of the $q$s. The extension of these results to
complex vector fields may be handled just as for complex scalar fields.

For the Dirac field of a non-Majorana spin $\tfrac12$ particle, Section 5.6 shows
that the anticommutator is
\[
\{\psi_\ell(x),\psi^\dagger_{\ell'}(y)\}
=\bigl[(-\gamma^\mu\partial_\mu+m)\beta\bigr]_{\ell\ell'}\Delta(x-y)
\]
and
\[
\{\psi_\ell(x),\psi_{\ell'}(y)\}=0.
\]
Here it would be inconsistent to take $\psi_\ell$ and $\psi_\ell^\dagger$ to be independent
canonical variables, because their anticommutator does not vanish at
equal times. It is conventional instead to define
\begin{equation}
q^\ell(x)\equiv\psi_\ell(x),
\tag{7.1.15}\label{eq:7.1.15}
\end{equation}
\begin{equation}
p_\ell(x)\equiv i\psi_\ell^\dagger(x).
\tag{7.1.16}\label{eq:7.1.16}
\end{equation}
It is easy then to see that \eqref{eq:7.1.15} and \eqref{eq:7.1.16} satisfy the canonical
anticommutation relations~\eqref{eq:7.1.1}--\eqref{eq:7.1.3}.

For any system of operators that satisfy commutation or anticommutation
relations like \eqref{eq:7.1.1}--\eqref{eq:7.1.3}, we may define a quantum mechanical
functional derivative: for an arbitrary bosonic functional $F[q(t),p(t)]$ of
$q^n(\mathbf{x},t)$ and $p_n(\mathbf{x},t)$ at a fixed time $t$, we define\footnote{We
are here using a notation that will be adopted henceforth; if $f(x,y)$ is a
function of two classes of variables collectively called $x$ and $y$, then
$F[f(y)]$ indicates a functional that depends on the values of $f(x,y)$ for
all $x$ at fixed $y$. By a bosonic functional we mean one in which each term
contains only even numbers of fermionic fields.}
\begin{equation}
\frac{\delta F[q(t),p(t)]}{\delta q^n(\mathbf{x},t)}
=i[p_n(\mathbf{x},t),F[q(t),p(t)]],
\tag{7.1.17}\label{eq:7.1.17}
\end{equation}
\begin{equation}
\frac{\delta F[q(t),p(t)]}{\delta p_n(\mathbf{x},t)}
=i[F[q(t),p(t)],q^n(\mathbf{x},t)].
\tag{7.1.18}\label{eq:7.1.18}
\end{equation}
This definition is motivated by the fact that if $F[q(t),p(t)]$ is written with
all $q$s to the left of all $p$s, then \eqref{eq:7.1.17} and \eqref{eq:7.1.18} are respectively just
the left- and right-derivatives with respect to $q^n$ and $p_n$. That is, for an
arbitrary c-number\footnote{Where $q^n$ and $p_n$ are bosonic or fermionic,
$\delta q^n$ and $\delta p_n$ are understood to commute or anticommute with all
fermionic operators, respectively, and to commute with all bosonic operators.}
variation $\delta q$ and $\delta p$ of the $q$s and $p$s, we have
\[
\delta F[q(t),p(t)]=\int d^3x\sum_n\left(
\delta q^n(\mathbf{x},t)\frac{\delta F[q(t),p(t)]}{\delta q^n(\mathbf{x},t)}
+\frac{\delta F[q(t),p(t)]}{\delta p_n(\mathbf{x},t)}\delta p_n(\mathbf{x},t)
\right).
\]
For more general functionals we need the definitions \eqref{eq:7.1.17} and \eqref{eq:7.1.18}
to pin down various signs and equal-time commutators that may appear.

In particular, $H_0$ is the generator of time-translations on free-particle
states in the sense that:
\begin{equation}
q^n(\mathbf{x},t)=\exp(iH_0t)q^n(\mathbf{x},0)\exp(-iH_0t),
\tag{7.1.19}\label{eq:7.1.19}
\end{equation}
\begin{equation}
p_n(\mathbf{x},t)=\exp(iH_0t)p_n(\mathbf{x},0)\exp(-iH_0t),
\tag{7.1.20}\label{eq:7.1.20}
\end{equation}
so the free-particle operators have the time-dependence
\begin{equation}
\dot q^n(\mathbf{x},t)=i[H_0,q^n(\mathbf{x},t)]
=\frac{\delta H_0}{\delta p_n(\mathbf{x},t)},
\tag{7.1.21}\label{eq:7.1.21}
\end{equation}
\begin{equation}
\dot p_n(\mathbf{x},t)=-i[p_n(\mathbf{x},t),H_0]
=-\frac{\delta H_0}{\delta q^n(\mathbf{x},t)}.
\tag{7.1.22}\label{eq:7.1.22}
\end{equation}
We recognize these as the familiar dynamical equations in the Hamiltonian
formalism.

The free-particle Hamiltonian is given as always by
\begin{equation}
H_0=\sum_{n,\sigma}\int d^3k\,
a^\dagger_{\mathbf{k},\sigma,n}a_{\mathbf{k},\sigma,n}
\sqrt{\mathbf{k}^2+m_n^2}.
\tag{7.1.23}\label{eq:7.1.23}
\end{equation}
This $H_0$ may be rewritten in terms of the $q$s and $p$s at time $t$. For instance,
it is easy to see that for a real scalar field, Eq.~\eqref{eq:7.1.23} is equal up to a
constant term to the functional
\begin{equation}
H_0=\int d^3x\left[\frac12p^2+\frac12(\boldsymbol{\nabla}q)^2
+\frac12m^2q^2\right].
\tag{7.1.24}\label{eq:7.1.24}
\end{equation}
To be more precise, using \eqref{eq:7.1.7} and the Fourier representation of the
scalar field $\phi$, we find that Eq.~\eqref{eq:7.1.24} becomes:
\begin{align}
H_0&=\frac12\int d^3k\,k^0\{a_{\mathbf{k}},a^\dagger_{\mathbf{k}}\}
\notag\\
&=\int d^3k\,k^0\left(a^\dagger_{\mathbf{k}}a_{\mathbf{k}}
+\frac12\delta^3(\mathbf{k}-\mathbf{k})\right).
\tag{7.1.25}\label{eq:7.1.25}
\end{align}
This is the same as Eq.~\eqref{eq:7.1.23}, except for the infinite constant term. Such
terms only affect the zero of energy, and have no physical significance in
the absence of gravity.\footnote{However, changes in such terms due to changes
in the boundary conditions for the fields, as for instance quantizing in the
space between parallel plates rather than in infinite space, are physically
significant, and have even been measured~\hyperref[ref:7.1]{[1]}.} Explicit forms for $H_0$ as a
functional of the $q$ and $p$ variables for other fields will be given in
Section 7.5.

It is usual in textbooks on quantum field theory to derive Eq.~\eqref{eq:7.1.25} as
a consequence of Eq.~\eqref{eq:7.1.24}, which in turn is derived from a Lagrangian
density. This seems to me backward, for Eq.~\eqref{eq:7.1.25} \textit{must} hold; if
some assumed free-particle Lagrangian did not give Eq.~\eqref{eq:7.1.25} up to
a constant term, we would conclude that it was the wrong Lagrangian.
Rather, we should ask what free-field Lagrangian gives Eq.~\eqref{eq:7.1.25} for
spinless particles, or more generally, gives the free-particle Hamiltonian
\eqref{eq:7.1.23}. This question may be answered by the well-known Legendre
transformation from the Hamiltonian to the Lagrangian; the free-field
Lagrangian is given by
\begin{equation}
L_0[q(t),\dot q(t)]\equiv
\sum_n\int d^3x\,p_n(\mathbf{x},t)\dot q^n(\mathbf{x},t)-H_0,
\tag{7.1.26}\label{eq:7.1.26}
\end{equation}
it being understood that $p_n$ is replaced everywhere by its expression in
terms of $q^n$ and $\dot q^n$ (and, as we shall see, perhaps some auxiliary fields
as well). For instance, from the Hamiltonian \eqref{eq:7.1.24} and \eqref{eq:7.1.7} we can
derive the free-field Lagrangian for a scalar field:
\begin{align}
L_0&=\int d^3x\left[p\dot q-\frac12p^2-\frac12(\boldsymbol{\nabla}q)^2
-\frac12m^2q^2\right]
\notag\\
&=\int d^3x\left[-\frac12\partial_\mu\phi\partial^\mu\phi
-\frac12m^2\phi^2\right].
\tag{7.1.27}\label{eq:7.1.27}
\end{align}
Whatever we suppose the complete Lagrangian of the scalar field may be,
this is the term that must be separated out and treated as a term of zeroth
order in perturbation theory. A similar exercise may be carried out for
the other canonical systems described in this section, but from now on we
shall content ourselves with guessing the form of the free-field Lagrangian
and then confirming that it gives the correct free-particle Hamiltonian.

We have seen that various free-field theories can be formulated in
canonical terms. It is then a short step to show that the same is true of the
interacting fields. We can introduce canonical variables in what is called
the `Heisenberg picture', defined by
\begin{equation}
Q^n(\mathbf{x},t)=\exp(iHt)q^n(\mathbf{x},0)\exp(-iHt),
\tag{7.1.28}\label{eq:7.1.28}
\end{equation}
\begin{equation}
P_n(\mathbf{x},t)=\exp(iHt)p_n(\mathbf{x},0)\exp(-iHt),
\tag{7.1.29}\label{eq:7.1.29}
\end{equation}
where $H$ is the full Hamiltonian. Because this is a similarity transformation
that commutes with $H$, the total Hamiltonian is the same functional of
the Heisenberg picture operators as it was of the $q$s and $p$s:
\[
H[Q,P]=e^{iHt}H[q,p]e^{-iHt}=H[q,p].
\]
Also, because Eqs.~\eqref{eq:7.1.28}--\eqref{eq:7.1.29} define a similarity transformation, the
Heisenberg picture operators again satisfy the canonical commutation or
anticommutation relations:
\begin{equation}
\begin{aligned}
[Q^n(\mathbf{x},t),P_{\bar n}(\mathbf{y},t)]
&=i\delta^3(\mathbf{x}-\mathbf{y})\delta^n_{\bar n},
&&\text{if either particle is a boson},\\
\{Q^n(\mathbf{x},t),P_{\bar n}(\mathbf{y},t)\}
&=i\delta^3(\mathbf{x}-\mathbf{y})\delta^n_{\bar n},
&&\text{if both particles are fermions}.
\end{aligned}
\tag{7.1.30}\label{eq:7.1.30}
\end{equation}
\begin{equation}
\begin{aligned}
[Q^n(\mathbf{x},t),Q^{\bar n}(\mathbf{y},t)]&=0,
&&\text{if either particle is a boson},\\
\{Q^n(\mathbf{x},t),Q^{\bar n}(\mathbf{y},t)\}&=0,
&&\text{if both particles are fermions}.
\end{aligned}
\tag{7.1.31}\label{eq:7.1.31}
\end{equation}
\begin{equation}
\begin{aligned}
[P_n(\mathbf{x},t),P_{\bar n}(\mathbf{y},t)]&=0,
&&\text{if either particle is a boson},\\
\{P_n(\mathbf{x},t),P_{\bar n}(\mathbf{y},t)\}&=0,
&&\text{if both particles are fermions}.
\end{aligned}
\tag{7.1.32}\label{eq:7.1.32}
\end{equation}
However, they now have the time-dependence
\begin{equation}
\dot Q^n(\mathbf{x},t)=i[H,Q^n(\mathbf{x},t)]
=\frac{\delta H}{\delta P_n(\mathbf{x},t)},
\tag{7.1.33}\label{eq:7.1.33}
\end{equation}
\begin{equation}
\dot P_n(\mathbf{x},t)=-i[P_n(\mathbf{x},t),H]
=-\frac{\delta H}{\delta Q^n(\mathbf{x},t)}.
\tag{7.1.34}\label{eq:7.1.34}
\end{equation}
For instance, we might take the Hamiltonian for a real scalar field as the
free-particle term \eqref{eq:7.1.24} plus the integral of a scalar interaction density
$\mathcal{H}_I$, so that in terms of Heisenberg-picture variables
\begin{equation}
H=\int d^3x\left[\frac12P^2+\frac12(\boldsymbol{\nabla}Q)^2
+\frac12m^2Q^2+\mathcal{H}_I(Q)\right].
\tag{7.1.35}\label{eq:7.1.35}
\end{equation}
In this case the canonical conjugate to $Q$ is given by the same formula as
for free fields:
\begin{equation}
P=\dot Q.
\tag{7.1.36}\label{eq:7.1.36}
\end{equation}
However, as we shall see, the relation between the canonical conjugates
$P_n(x)$ and the field variables and their time-derivatives is in general not
the same as for the free-particle operators, but must be inferred from
Eqs.~\eqref{eq:7.1.33} and \eqref{eq:7.1.34}.
